\chapter{Experimental Methodology}

\section{Evaluation philosophy}

The thesis evaluation strategy is incremental:

\begin{enumerate}
    \item validate solver behavior in controlled simulation,
    \item validate software integration under runtime constraints,
    \item validate inverse estimation quality under increasing uncertainty.
\end{enumerate}

This staged methodology prevents mixing algorithmic and systems failures in a
single experiment.

\section{Evaluation dimensions}

Three dimensions are measured across work packages:

\begin{itemize}
    \item \textbf{efficiency}: solve time and runtime overhead,
    \item \textbf{behavior}: safety/goal/lane metrics and qualitative
    trajectory signatures,
    \item \textbf{estimation}: preference-order identification quality and
    downstream prediction impact.
\end{itemize}

\section{Forward-solver benchmarks}

For forward solving, the baseline metrics from \cite{delasheras:2025} are reused
and extended:

\begin{itemize}
    \item solve time per receding-horizon step,
    \item compile/setup overhead,
    \item convergence or stabilization behavior versus IBR rounds.
\end{itemize}

\section{Integration benchmarks}

For \texttt{patos}, planned metrics include:

\begin{itemize}
    \item manager cycle consistency,
    \item solver timeout/fallback rate,
    \item command tracking and safety-clearance indicators,
    \item behavior under multi-robot load.
\end{itemize}

\section{Inverse-estimation benchmarks}

For \texttt{InverseGoopSolver}, planned metrics include:

\begin{itemize}
    \item exact-match rate of inferred order,
    \item rank-distance metrics between inferred and true order,
    \item trajectory mismatch before/after estimation,
    \item robustness against observation noise and partial state access.
\end{itemize}

\begin{figure}[tbp]
    \centering
    \fbox{\parbox{0.92\linewidth}{
    \textbf{Placeholder figure: Metric dashboard layout.}
    Panel view with efficiency, safety/behavior, and estimation metrics to be
    reported uniformly across experiments.
    }}
    \caption{Planned unified metric dashboard for final results chapters.}
    \label{fig:metric_dashboard}
\end{figure}

\section{Reproducibility protocol}

To keep comparisons reliable, each final experiment set will record:

\begin{itemize}
    \item exact configuration parameters,
    \item software commit references,
    \item random seeds,
    \item generated plots and serialized summary artifacts.
\end{itemize}

This protocol will be used for both simulation runs and embodied Duckietown
experiments.
